% !TeX program = pdflatex
% Kandydacki master DT. Kompilować z books/pl, po functional_analysis.tex.
\documentclass[11pt,a4paper]{article}
\usepackage[T1]{fontenc}
\usepackage[utf8]{inputenc}
\usepackage[polish]{babel}
\usepackage{lmodern}
\usepackage{amsmath,amssymb,amsthm,mathtools}
\usepackage{microtype}
\usepackage{geometry}
\usepackage{enumitem}
\usepackage{tocloft}
\usepackage{xcolor}
\usepackage{tikz}
\usetikzlibrary{arrows.meta}
\usepackage{xr-hyper}
\usepackage[unicode]{hyperref}
\geometry{margin=2.5cm}
\hypersetup{colorlinks=true,linkcolor=blue!45!black,urlcolor=blue!45!black,pdftitle={Teoria dystrybucji i analiza harmoniczna}}
\setlength{\parindent}{0pt}
\setlength{\parskip}{4pt}
\setlength{\emergencystretch}{2em}
\setlength{\cftsecnumwidth}{3em}
\setlength{\cftsubsecnumwidth}{4.5em}
% Moduły 15A/15B mają numery podpodsekcji w rodzaju 15B.12.8.
\setlength{\cftsubsubsecnumwidth}{6em}
\setlist[enumerate]{topsep=4pt,itemsep=2pt}
\setlist[itemize]{topsep=4pt,itemsep=2pt}
\numberwithin{equation}{section}
\theoremstyle{plain}
\newtheorem{twierdzenie}{Twierdzenie}[section]
\newtheorem{stwierdzenie}[twierdzenie]{Stwierdzenie}
\newtheorem{lemat}[twierdzenie]{Lemat}
\newtheorem{wniosek}[twierdzenie]{Wniosek}
\theoremstyle{definition}
\newtheorem{definicja}[twierdzenie]{Definicja}
\newtheorem{przyklad}[twierdzenie]{Przykład}
\theoremstyle{remark}
\newtheorem{uwaga}[twierdzenie]{Uwaga}
\newcommand{\R}{\mathbb R}
\newcommand{\C}{\mathbb C}
\newcommand{\N}{\mathbb N}
\newcommand{\K}{\mathbb K}
\newcommand{\cl}[1]{\overline{#1}}
\newcommand{\TVS}{\mathrm{TVS}}
\DeclareMathOperator{\lin}{lin}
\DeclareMathOperator{\supp}{supp}
\DeclareMathOperator{\id}{id}
\title{Teoria dystrybucji i analiza harmoniczna}
\author{}
\date{Kandydat restrukturyzacji -- 28 listopada 2026}

% BEGIN GLOBAL CODE LISTINGS (generated from latex/pl/code-listings.tex)
\usepackage{xcolor}
\usepackage{listings}
\definecolor{codebg}{RGB}{248,249,251}
\definecolor{codeborder}{RGB}{180,191,203}
\definecolor{codecomment}{RGB}{61,113,79}
\definecolor{codekeyword}{RGB}{41,78,132}
\lstdefinelanguage{Julia}{
  morekeywords={abstract,baremodule,begin,break,catch,const,continue,do,else,elseif,
    end,export,false,finally,for,function,global,if,import,in,isa,let,local,macro,
    module,mutable,primitive,quote,return,struct,true,try,type,using,where,while,
    missing,nothing,println,print},
  sensitive=true,morecomment=[l]{\#},morecomment=[n]{\#=}{=\#},
  morestring=[b]",
}
\lstdefinestyle{skrypt}{
  backgroundcolor=\color{codebg},frame=single,rulecolor=\color{codeborder},
  framesep=5pt,frameround=tttt,framerule=.5pt,
  basicstyle=\ttfamily\footnotesize,keywordstyle=\color{codekeyword}\bfseries,
  commentstyle=\color{codecomment}\itshape,stringstyle=\color{purple!65!black},
  numbers=left,numberstyle=\tiny\color{gray!75!black},numbersep=8pt,
  stepnumber=1,breaklines=true,breakatwhitespace=false,
  columns=fullflexible,keepspaces=true,showstringspaces=false,tabsize=4,
  linewidth=.90\linewidth,captionpos=b,aboveskip=11pt,belowskip=12pt,
  xleftmargin=12pt,xrightmargin=2pt,
}
\lstdefinestyle{pythonpde}{style=skrypt,language=Python}
\lstdefinestyle{python}{style=skrypt,language=Python}
\lstdefinestyle{julia}{style=skrypt,language=Julia}
\lstset{style=skrypt}
\renewcommand{\lstlistingname}{Kod}
\renewcommand{\lstlistlistingname}{Spis kodów}
% END GLOBAL CODE LISTINGS

% AF pozostaje odrębnym źródłem trzech wcześniejszych wyników.
% Kompilacja: najpierw functional_analysis.tex, następnie ten skrypt.
% Prefiks zapobiega kolizji z 211 etykietami własnych rozdziałów.
\externaldocument[AF-ext-]{functional_analysis}
\makeatletter
\newcommand{\importFAreference}[1]{%
  \@ifundefined{r@AF-ext-#1}{%
    \PackageWarningNoLine{distribution-theory}{Missing AF reference #1; compile functional_analysis first}%
  }{\expandafter\let\csname r@#1\expandafter\endcsname\csname r@AF-ext-#1\endcsname}}
\importFAreference{cor:frechet-barrelled}
\importFAreference{thm:banach-steinhaus-lcs}
\importFAreference{thm:point-closed-subspace}
\makeatother
\providecommand{\atlasref}[2]{#2}
\begin{document}
\maketitle
\noindent Wersja kandydacka: numeracja docelowa, pominięto rozdziały bez istniejącego materiału.\par
\tableofcontents
\clearpage
\setcounter{part}{0}
\part{Podstawy teorii dystrybucji}
\setcounter{section}{0}
\input{../../shared/pl/distribution_theory/restructure/01-przestrzenie-funkcji-gladkich-i-schwartza.tex}
\setcounter{section}{1}
\input{../../shared/pl/distribution_theory/restructure/02-funkcje-testowe-i-scisle-granice-induktywne.tex}
\setcounter{section}{2}
\input{../../shared/pl/distribution_theory/restructure/03-przestrzen-dystrybucji-i-topologie.tex}
\setcounter{section}{3}
\input{../../shared/pl/distribution_theory/restructure/04-lokalnosc-nosniki-i-przedluzanie.tex}
\setcounter{section}{4}
\input{../../shared/pl/distribution_theory/restructure/05-pochodne-dystrybucyjne-i-transpozycja-operatorow.tex}
\setcounter{section}{5}
\input{../../shared/pl/distribution_theory/restructure/06-dystrybucje-o-zwartym-nosniku.tex}
\setcounter{section}{6}
\input{../../shared/pl/distribution_theory/restructure/07-splot-dystrybucji-i-regularyzacja.tex}
\setcounter{section}{7}
\input{../../shared/pl/distribution_theory/restructure/08-skoki-powierzchniowe-i-iloczyny-dystrybucji.tex}
\setcounter{part}{1}
\part{Transformacja Fouriera}
\setcounter{section}{8}
\input{../../shared/pl/distribution_theory/restructure/09-transformacja-fouriera-schwartz-l1.tex}
\setcounter{section}{9}
\input{../../shared/pl/distribution_theory/restructure/10-dystrybucje-temperowane-i-fourier.tex}
\setcounter{section}{10}
\input{../../shared/pl/distribution_theory/restructure/11-plancherel-i-zgodnosc-transformat.tex}
\setcounter{section}{11}
\input{../../shared/pl/distribution_theory/restructure/12-splot-funkcji-i-young.tex}
\setcounter{section}{12}
\input{../../shared/pl/distribution_theory/restructure/13-interpolacja-zespolona-i-hausdorff-young.tex}
\setcounter{section}{13}
\input{../../shared/pl/distribution_theory/restructure/14-transformacja-fouriera-miar.tex}
\setcounter{section}{14}
\input{../../shared/pl/distribution_theory/restructure/15-transformacja-fouriera-funkcji-radialnych.tex}
\setcounter{section}{15}
\input{../../shared/pl/distribution_theory/restructure/16-transformata-dystrybucji-o-zwartym-nosniku.tex}
\setcounter{part}{2}
\part{Struktura topologiczna i dalsza teoria dystrybucji}
\setcounter{section}{16}
\input{../../shared/pl/distribution_theory/restructure/17-montel-i-nuklearnosc.tex}
\setcounter{section}{17}
\input{../../shared/pl/distribution_theory/restructure/18-topologiczne-iloczyny-tensorowe.tex}
\setcounter{section}{18}
\input{../../shared/pl/distribution_theory/restructure/19-twierdzenie-o-jadrze-schwartza.tex}
\setcounter{section}{19}
\input{../../shared/pl/distribution_theory/restructure/20-induktywne-iloczyny-tensorowe.tex}
\setcounter{section}{20}
\input{../../shared/pl/distribution_theory/restructure/21-rozwiazania-fundamentalne.tex}
\setcounter{part}{3}
\part{Szeregi Fouriera}
\setcounter{section}{23}
\input{../../shared/pl/distribution_theory/restructure/24-uklady-ortogonalne-i-wspolczynniki-fouriera.tex}
\setcounter{section}{25}
\input{../../shared/pl/distribution_theory/restructure/26-jadra-fejera-i-poissona.tex}
\setcounter{section}{26}
\input{../../shared/pl/distribution_theory/restructure/27-zbieznosc-jednostajna-i-punktowa.tex}
\setcounter{part}{6}
\part{Całki osobliwe i klasyczna analiza harmoniczna}
\setcounter{section}{40}
\input{../../shared/pl/distribution_theory/restructure/41-operator-maksymalny.tex}
\setcounter{section}{41}
\input{../../shared/pl/distribution_theory/restructure/42-interpolacja-rzeczywista.tex}
\setcounter{section}{42}
\input{../../shared/pl/distribution_theory/restructure/43-wartosci-glowne-i-mnozniki-l2.tex}
\setcounter{section}{45}
\input{../../shared/pl/distribution_theory/restructure/46-rozklad-i-teoria-calderona-zygmunda.tex}
\setcounter{section}{49}
\input{../../shared/pl/distribution_theory/restructure/50-potencjaly-bessela-sobolew-i-nieoznaczonosc.tex}
\setcounter{part}{7}
\part{Zaawansowana analiza harmoniczna}
\setcounter{section}{56}
\input{../../shared/pl/distribution_theory/restructure/57-splot-na-grupach.tex}
\end{document}
